\section{News and compatibility}\label{news}

\subsection{With version 5.15}
\begin{itemize}
  \item \tkzcname{tkzInterLC} with a line passing through the center of the circle: the order of the two points no longer depends on rounding. \tkzname{tkzSecondPoint} is the point located in the direction of the line (from its first point to its second point); the results are the same with and without the option \tkzname{lua}. See section \ref{interLC-center}.
\end{itemize}

\subsection{With version 5.14}
\begingroup\sloppy

This version only corrects bugs; the syntax does not change and no macro has been added.

\begin{itemize}
  \item \tkzcname{tkzDefCircle} with the options \tkzname{orthogonal from} or \tkzname{orthogonal through} no longer disturbs the following calls of \tkzcname{tkzDefCircle}.
  \item The options of the following macros now only apply to the current call: \tkzcname{tkzShowLine}, \tkzcname{tkzDefSimilitudeCenter}, \tkzcname{tkzCalcLength}, \tkzcname{tkzGrid}, \tkzcname{tkzClip}, \tkzcname{tkzDefEquiPoints}, \tkzcname{tkzLegend}.
  \item \tkzcname{tkzIsLinear} gave a wrong answer (always true); \tkzcname{tkzIsLinear} and \tkzcname{tkzIsOrtho} now use a tolerance that depends on the lengths.
  \item The result of \tkzcname{tkzPowerCircle} no longer changes after another calculation.
  \item The inversion of a circle (\tkzcname{tkzDefCircleBy}) gives a line only if the circle passes through the center of inversion.
  \item \tkzcname{tkzDefCircle}\code{[spieker]} gives the radius in cm, like the other options.
  \item The Feuerbach point and the Feuerbach triangle, obtained with the option \tkzname{feuerbach} of \tkzcname{tkzDefTriangleCenter} and \tkzcname{tkzDefSpcTriangle}, are now exact.
  \item \tkzcname{tkzDrawPolySeg}: the arrow at the start of the path has the right direction.
  \item \tkzcname{tkzCompass} accepts a decimal value for \tkzname{delta}; \tkzcname{tkzProtractor}\code{[return]} works; \tkzcname{tkzLabelRegPolygon} uses its option \tkzname{sep}; \tkzcname{tkzDrawX} uses its option \tkzname{step}.
  \item The old syntax of version 3, for example |\tkzDrawCircle[euler](A,B,C)|, now gives an error message instead of an endless compilation.
  \item The aliases \tkzcname{tkzDefNinePointsCircle}, \tkzcname{tkzFeuerBachCircle} and \tkzcname{tkzIncentralTriangle} work again.
\end{itemize}

Two changes may modify your figures:
\begin{itemize}
  \item \tkzcname{tkzSetUpCompass}\code{[delta=...]} only acts on the compass marks. To extend the arcs drawn by \tkzcname{tkzDrawArc}, \tkzcname{tkzMarkArc} and \tkzcname{tkzLabelArc}, use \tkzcname{tkzSetUpArc}\code{[delta=...]}. Before this version, the same value was used for both.
  \item The \TIKZ\ style \tkzname{help lines} is no longer modified by the package.
\end{itemize}

\endgroup

\subsection{With version 5.10}
\begin{itemize}
  \item Added a French documentation.
  \item Added the \code{mini} option. You can use this option with the \tkzNamePack{tkz-elements} package. Only the modules required for drawing will be loaded. This option is currently only useful if you are using \tkzNamePack{tkz-elements}.
\end{itemize}

\subsection{From version 3 to version 5}

\tkzname{\tkznameofpack} 4 was a simplified version of its predecessor, independent of \tkzname{tkz-base}. Most macros of \tkzname{tkz-euclide} 3 have been retained, but the special circles are now defined with \tkzcname{tkzDefCircle} and then drawn with \tkzcname{tkzDrawCircle}. The unit is now the \tkzname{cm}. If you need some macros from \tkzname{tkz-base}, you may need to use \tkzcname{tkzInit}.

Version 5 does not change anything for the user. With the option \code{lua} (|\usepackage[lua]{tkz-euclide}|), the calculations are done by Lua: you must then compile with LuaLaTeX. See section \ref{calc_with_lua}.



\endinput